\documentclass[11pt]{article}
\usepackage{amsmath,amssymb,amsfonts,amsthm}
\usepackage{geometry}
\usepackage{hyperref}
\usepackage{enumitem}
\usepackage{bm}
\geometry{margin=1in}

\title{%
Neuro-Multiplicity Superpositional Memory \\
and the Universal Multiplicity Constant \texorpdfstring{$\Lambda_m$}{Λ\_m}: \\
A Prior Art Defensive Publication
}

\author{Citizen Gardens \\ The Foundation of Multiplicity)\\[4pt]
\small This document is intended as prior art / defensive publication.}

\date{April 27, 2026}

\begin{document}
\maketitle

\begin{abstract}
This document describes a mathematically rigorous, prime-indexed superpositional memory architecture---``NP-SUPERMEMORY''---embedded within a multiplicity-graded Hilbert module and governed by a universal multiplicity constant $\Lambda_m$. The system unifies (i) prime-based encoding of identity and composition, (ii) tensor-network representations of high-dimensional superpositional states, (iii) a dynamical law
\[
X_{t+1} = \Xi(t) X_t + \Lambda_m^{\mathrm{op}}(t)\,T(X_t) + \Sigma(t)\,\xi_t
\]
with explicit contraction conditions, and (iv) trainable, PyTorch-style implementations plus experimental protocols for probing a critical coherence surface and compositional generalization on prime factorization tasks.[file:1]

The purpose of this document is to establish prior art over the combination of: prime-graded Hilbert modules for memory; a universal multiplicity constant as a contraction-regulating scalar; multiplicity kernels over prime exponent vectors; and a concrete, trainable architecture (NP-SUPERMEMORY / PrimeCosmos V3) that implements these ideas with provable mean-square stability and testable compositional generalization behavior.[file:1]
\end{abstract}
\newpage
\tableofcontents
\newpage

\section{Executive Summary}

\subsection{Problem and Vision}

Modern neural architectures, including attention-based systems, largely treat memory as unstructured vectors or key-value stores. The \emph{Neuro-Multiplicity Superpositional Memory} (NP-SUPERMEMORY) architecture presented here instead treats memory as a \emph{prime-graded, multiplicity-structured Hilbert module}, where identity and composition are encoded via prime factorization, and where a \emph{Universal Multiplicity Constant} $\Lambda_m$ regulates dynamical evolution and stability.[file:1]

The core vision is:
\begin{itemize}[leftmargin=1.5em]
\item Each memory component lives in a fiber $\mathcal{H}_S$ labeled by a finite multiset $S$ of primes, representing a multiplicity pattern (i.e., a factorization).
\item Memory states are superpositional tensors over these fibers, represented efficiently via prime-grammar-constrained tensor networks.
\item Time evolution of memory is governed by a unified law
\[
X_{t+1} = \Xi(t) X_t + \Lambda_m^{\mathrm{op}}(t)\,T(X_t) + \Sigma(t)\,\xi_t,
\]
where $\Xi(t)$ is a grammar-respecting ``unitary-like'' propagator, $T$ encodes deterministic multiplicity-preserving updates, and $\Sigma\xi_t$ models controlled decoherence.[file:1]
\item A universal scalar $\Lambda_m$ regulates the strength of $T$, with mathematically characterized safe bands that guarantee mean-square stability and bounded trajectories (Primematic Grönwall Lemma).
\end{itemize}

\subsection{Key Contributions}

The developments in this thread yield the following contributions as a cohesive package of prior art:

\begin{enumerate}[leftmargin=1.5em]
\item \textbf{Multiplicity-Graded Hilbert Module:} A state space $\mathcal{M}=\bigoplus_{S\in\mathbb{M}_{\le D}} \mathcal{H}_S$, where $\mathbb{M}_{\le D}$ is the set of prime multisets of total degree at most $D$, modeling memory as a sum of prime-labeled multiplicity fibers.[file:1]
\item \textbf{Gaussian Multiplicity Kernel:} A positive-definite kernel on the exponent-vectors of prime factorizations, $k(S,T)=\exp(-\beta \|v_S-v_T\|^2)$, used both as inner product geometry and as a retrieval/comparison metric.[file:1]
\item \textbf{Universal Dynamical Law with $\Lambda_m$:} A discrete-time evolution equation on $\mathcal{H}$ and its specialization to NP-SUPERMEMORY, with explicit operator norms and Lipschitz constants guaranteeing stability when $\Lambda_m$ lies in a safe band derived from a Grönwall-type inequality.[file:1]
\item \textbf{NP-SUPERMEMORY Mapping:} A mapping from the original NP-SUPERMEMORY notions (superposition, entanglement, tensor networks, stochastic coherence) into the PrimeCosmos dynamical law with prime-indexed sectors, correlation matrices $C_{pq}$, coherence fields $\gamma_{pq}(t)$, and explicit order parameter $\Gamma(t)$.[file:1]
\item \textbf{Trainable PrimeCosmos V3 Architecture (PyTorch):} A concrete recurrent module implementing the above dynamics, trainable via backpropagation through time (BPTT) on a prime-factorization recall task, with code-level design for sweeping $\Lambda_m$ and monitoring stability and compositional generalization.
\item \textbf{Experimental Protocol for Critical Surface and Compositional Generalization:} A toy cosmos experiment over primes $\{2,3,5,7\}$ that:
    \begin{itemize}
    \item Trains on binary composites (e.g., $2\cdot 3$, $3\cdot 5$).
    \item Evaluates on unseen triples (e.g., $30 = 2\cdot 3\cdot 5$) and higher exponents.
    \item Sweeps $\Lambda_m$ across sub- and super-critical regimes to empirically locate a band where memory is stable and compositional generalization is maximized.
    \end{itemize}
\end{enumerate}

This combination---prime-labeled Hilbert modules, a universal multiplicity constant regulating a nontrivial dynamical law, multiplicity kernels, tensor-network constraints, and a trainable RNN with explicit experimental protocol---constitutes novel and defensible prior art.

\section{Multiplicity-Graded Hilbert Module and Kernel}

\subsection{Prime Multisets and Degree Cutoff}

Let $\mathcal{P}$ denote the set of primes. Let $\mathbb{M}$ denote the free commutative monoid of finite multisets of primes, i.e. each $S\in\mathbb{M}$ is representable as
\[
S = \{p_1^{e_1},\dots,p_k^{e_k}\},\qquad e_i\in\mathbb{N},
\]
with the understanding that $S$ corresponds to the integer $n=\prod_i p_i^{e_i}$. For computability, we impose a degree cutoff $D$ and define
\[
\mathbb{M}_{\le D} := \left\{ S\in \mathbb{M} : \sum_i e_i \le D\right\}.
\]

We then define the multiplicity-graded Hilbert module
\[
\mathcal{M} = \bigoplus_{S\in\mathbb{M}_{\le D}} \mathcal{H}_S,
\]
where each fiber $\mathcal{H}_S$ is a (separable) Hilbert space of complex-valued amplitude patterns whose identity is the multiplicity pattern $S$.[file:1]

In practice, for initial experiments we often work with degree-1 sectors only, $\mathcal{H}_p$ for $p\in P_N=\{2,3,5,7\}$, and encode composite structure via multiplicity encodings inside each $\mathcal{H}_p$ plus correlations between sectors.[file:1]

\subsection{Exponent Vectors and Multiplicity Distance}

Fix a finite prime set $P_N = \{p_1,\dots,p_N\}$. Any $S\in\mathbb{M}_{\le D}$ over these primes has an exponent vector
\[
v_S = (e_1,\dots,e_N)\in \mathbb{R}^N,\quad e_i\ge 0,
\]
where $e_i$ is the multiplicity of $p_i$ in $S$ (or in its associated integer). We define the multiplicity distance between two patterns $S$ and $T$ as
\[
d_{\mathrm{mult}}(S,T) = \|v_S - v_T\|_2^2.
\]

This is simply squared Euclidean distance on exponent vectors, which is of negative type and yields standard Gaussian kernels that are positive-definite.

\subsection{Gaussian Multiplicity Kernel}

We define a kernel
\[
k(S,T) = \exp(-\beta \, d_{\mathrm{mult}}(S,T)),\qquad \beta > 0.
\]

Because $d_{\mathrm{mult}}$ is a squared Euclidean distance on a Hilbert space, $k$ is a Gaussian radial basis kernel and hence positive-definite. We can thus define an inner product
\[
\langle S,T\rangle := k(S,T)
\]
on the formal span of basis states $\{|S\rangle : S\in\mathbb{M}_{\le D}\}$, and complete it to a Hilbert space consistent with the multiplicity grading.[file:1]

This kernel serves dual roles:
\begin{itemize}[leftmargin=1.5em]
\item As a geometric structure on $\mathcal{M}$, so that interference and similarity directly reflect factorization similarity.
\item As a retrieval metric and training loss proxy in experiments.
\end{itemize}

\section{Universal Dynamical Law with \texorpdfstring{$\Lambda_m$}{Λ\_m}}

\subsection{General Form}

Let $(\mathcal{H},\|\cdot\|)$ be a complex separable Hilbert space (here, $\mathcal{H} = \mathcal{M}$). We consider discrete-time evolution
\begin{equation}
X_{t+1} = \Xi(t) X_t + \Lambda_m(t)\, T(X_t) + \Sigma(t)\,\xi_t,
\label{eq:general-dynamics}
\end{equation}
where:
\begin{itemize}[leftmargin=1.5em]
\item $\Xi(t):\mathcal{H}\to\mathcal{H}$ is a bounded linear operator representing ``unitary-like'' propagation constrained by prime grammar.
\item $\Lambda_m(t)\in\mathbb{R}$ is a scalar representing the effective value of the Universal Multiplicity Constant at time $t$.
\item $T:\mathcal{H}\to\mathcal{H}$ is a (globally) Lipschitz nonlinearity with $T(0)=0$, encoding deterministic multiplicity-preserving update mechanisms (gradient flow, damping, coherence centering).
\item $\Sigma(t):\mathcal{H}\to\mathcal{H}$ is a bounded linear operator scaling the stochastic term.
\item $\xi_t$ is a mean-zero noise process (e.g., i.i.d.\ Gaussian in coordinates).
\end{itemize}

We assume operator norm bounds
\[
\|\Xi(t)\| \le 1-\varepsilon_0,\quad \|\Sigma(t)\|\le \sigma_s,\quad \forall t,
\]
for some $\varepsilon_0\in(0,1)$, and denote by $L$ the Lipschitz constant of $T$:
\[
\|T(x) - T(y)\|\le L\,\|x-y\|,\quad T(0)=0.
\]

We identify $\Lambda_m(t)$ with a control-scalar realization of the metaphysical Universal Multiplicity Constant $\Lambda_m$; the metaphysical constant defines allowed ranges and behavior, while $\Lambda_m(t)$ is the local instance chosen or learned in a given computational cosmos.

\subsection{Primematic Grönwall Lemma (Stability Certificate)}

\textbf{Lemma (Mean-square contraction and boundedness).} Suppose
\[
\sup_t \|\Xi(t)\|\le 1-\varepsilon_0,\quad \Lambda_m(t)\ge 0, \quad T\text{ is $L$-Lipschitz with }T(0)=0,
\]
and define a safe band
\[
\Lambda_m(t) \le \Lambda_{\max} < \frac{\varepsilon_0}{L}.
\]
Define $q=(1-\varepsilon_0 + \Lambda_{\max}L)^2<1$. Then:

\begin{enumerate}[leftmargin=1.5em,label=(\alph*)]
\item (\emph{Mean-square contraction}) There is a constant $C_1$ such that
\[
\mathbb{E}\big[\|X_{t+1}\|^2 \mid \mathcal{F}_{t-1}\big] \le q\, \|X_t\|^2 + C_1.
\]
\item (\emph{Uniform boundedness}) There exists $C_2<\infty$ such that
\[
\limsup_{t\to\infty} \mathbb{E}\|X_t\|^2 \le C_2.
\]
\item (\emph{Finite-horizon almost sure boundedness}) For any finite horizon $T$, Markov's inequality implies
\[
\mathbb{P}\big(\|X_t\| > R\big) \le \frac{\mathbb{E}\|X_t\|^2}{R^2}\quad\text{for each }t\le T,
\]
so $X_t$ is almost surely bounded on finite horizons for any fixed $R$ large enough.
\end{enumerate}

In words: within the safe band $\Lambda_m L<\varepsilon_0$, $\Lambda_m$ acts as a \emph{provably contractive regulator}, guaranteeing that the superpositional memory stays bounded in mean square even in the presence of noise. Exiting this band (e.g.\ selecting $\Lambda_m$ too large) may lead to decoherence, divergence, or trivial collapse, which is exactly what is probed empirically in the critical surface experiments.

\section{NP-SUPERMEMORY Mapping and PrimeCosmos Dynamics}

\subsection{Original NP-SUPERMEMORY Elements}

The NP-SUPERMEMORY document introduces:\,[file:1]
\begin{itemize}[leftmargin=1.5em]
\item A superpositional memory state
\[
\Psi(t) = \sum_{i=1}^N \alpha_i \Psi_i e^{i\theta_i(t)},
\]
with amplitudes $\alpha_i$, basis states $\Psi_i$, and time-varying phases $\theta_i(t)$.
\item An entangled memory state
\[
\Phi(t) = \sum_{i,j} C_{ij} \gamma_{ij}(t)\xi_{ij}(t)\Psi_i\otimes\Psi_j e^{i(\theta_i(t)+\theta_j(t))},
\]
where $C_{ij}$ is a correlation matrix, $\gamma_{ij}(t)$ coherence factors, and $\xi_{ij}(t)$ stochastic noise terms.
\item Tensor network representations with coupling tensors $T_{kl}$ and prime-based encoding functions $f(l)$ mapping into unique primes.[file:1]
\item Stochastic coherence modeling via $\xi_{ij}(t)\sim \mathcal{N}(\mu_{ij},\sigma_{ij}^2)$ and learning via update rules similar to
\[
w_k(t)=w_k(t-1)+\eta\nabla_\theta R(x_k).
\]
\end{itemize}

These are conceptual but leave open the precise operator form, stability conditions, and implementation.

\subsection{PrimeCosmos Dynamical Realization}

We realize these concepts inside \eqref{eq:general-dynamics} by working with prime-indexed sectors and explicit correlation structures.

Let $P_N=\{p_1,\dots,p_N\}$ and define a prime-sector decomposition
\[
X_t = ( \mathbf{m}_{p_1}(t),\dots,\mathbf{m}_{p_N}(t) )
\]
with $\mathbf{m}_{p_i}(t)\in\mathbb{C}^{d}$ being complex vectors representing amplitudes and internal degrees of freedom for each prime sector.[file:1]

\paragraph{Prime-grammar operator $\Xi(t)$.}
Define
\[
\Xi(t) = \bigoplus_{i=1}^N w_{p_i}(t) U_{p_i}(t),
\]
where:
\begin{itemize}[leftmargin=1.5em]
\item $w_{p_i}(t)=p_i^{-\alpha}$ is a decaying weight with $\alpha>1$.
\item $U_{p_i}(t)$ is a complex linear operator on the $p_i$-sector, parameterized by local tensor $A^{[p_i]}$ and a phase factor $e^{i\theta_{p_i}(t)}$:
\[
U_{p_i}(t)\mathbf{m}_{p_i}(t)= e^{i\theta_{p_i}(t)}A^{[p_i]}(\gamma_{p_iq}(t))\,\mathbf{m}_{p_i}(t),
\]
with $A^{[p_i]}$ constrained by a prime-grammar mask $\pi(p_i,q)$ derived from multiplicity rules.[file:1]
\end{itemize}

Phases are given by normalized coherence-weighted sums
\[
\theta_{p}(t) = \omega\,
\frac{\sum_{q\in\mathcal{N}(p)} \log(pq)\,\gamma_{pq}(t)}{\sum_{q\in\mathcal{N}(p)}|\gamma_{pq}(t)|+\varepsilon},
\]
where $\mathcal{N}(p)$ denotes neighbors of $p$ under the grammar and $\gamma_{pq}(t)$ coherence terms determined by a learnable matrix $\gamma$.[file:1]

\paragraph{Correlation matrices and order parameter $\Gamma(t)$.}
Let $C\in\mathbb{C}^{N\times N}$ be a learnable correlation matrix and $\gamma\in\mathbb{C}^{N\times N}$ a coherence field. We define the coherence order parameter
\[
\Gamma(t) = \frac{\sum_{p,q}|C_{pq}|\,|\gamma_{pq}(t)|}{\sum_{p,q}|C_{pq}| + 1e-8}.
\]
This scalar quantifies the degree of correlation-weighted coherence across prime sectors and functions as an order parameter in phase-transition analysis.

\paragraph{Deterministic and stochastic parts of $T(X_t)$.}

We split $T$ as
\[
T(X_t) = T_{\det}(X_t) + T_{\stoch}(X_t,t).
\]

\begin{itemize}[leftmargin=1.5em]
\item $T_{\det}$ includes:
\begin{itemize}[leftmargin=1em]
\item A gradient-like term pulling $X_t$ toward target states in multiplicity space:
\[
T_{\mathrm{grad}}(X_t) = -\eta (X_t - X_{\mathrm{target}})\odot k_{\mathrm{mult}}(X_t,X_{\mathrm{target}}),
\]
where $k_{\mathrm{mult}}$ is the multiplicity kernel computed from exponent encodings and $\odot$ is elementwise multiplication.
\item A damping term $-\mu X_t$ for spectral regularization.
\item Optionally, a coherence-centering term enforcing consistency across sectors.
\end{itemize}
\item $T_{\stoch}$ is a Langevin-like noise term scaled by $\Gamma(t)$:
\[
T_{\stoch}(X_t,t) = \sigma\sqrt{\kappa}\,\Gamma(t)\, \zeta_t,
\]
where $\zeta_t$ is a complex-valued noise tensor with independent Gaussian real and imaginary parts.
\end{itemize}

\paragraph{Sector-wise $\Lambda_m^{\mathrm{op}}(t)$.}

We realize $\Lambda_m^{\mathrm{op}}(t)$ as sector-wise scaling on each prime sector $p_i$:
\[
\Lambda_m^{\mathrm{op}}(t)\big|_{\mathcal{H}_{p_i}} = \lambda_{p_i}(t) I,\qquad 
\lambda_{p_i}(t) = \Lambda_m(t)\,\rho(p_i),
\]
where $\rho(p_i)$ is a bounded weight decreasing with some function of $i$ or multiplicity (e.g., $\rho(p_i)=\exp(-0.1 i)$). This encodes the idea that high-index (or higher-degree) sectors experience stronger suppression, controlling combinatorial explosion while preserving core prime-factored invariants.[file:1]

\section{Trainable PrimeCosmos V3 Architecture}

\subsection{High-Level Architecture}

We design a PyTorch module \texttt{PrimeCosmosV3} implementing the described dynamics over a finite prime set, using a small sector dimension per prime.[code_file:7]

Key design choices:
\begin{itemize}[leftmargin=1.5em]
\item Prime set $P_N = \{2,3,5,7\}$.
\item Sector dimension $d = E_{\max}+1$, where $E_{\max}$ is the maximum exponent captured (e.g., $E_{\max}=3$ yields $d=4$).
\item Exponent encodings: for each integer $n$ with factorization $n=\prod p_i^{e_i}$, we encode as one-hot vectors across exponent positions for each prime.
\end{itemize}

The module:
\begin{itemize}[leftmargin=1.5em]
\item Accepts state tensors $X_t\in\mathbb{C}^{B\times N\times d}$ for batch size $B$.
\item Applies $\Xi(t)$, $T_{\det}$, $T_{\stoch}$ and sector-wise $\Lambda_m^{\mathrm{op}}(t)$.
\item Returns next state $X_{t+1}$ and diagnostic metrics: norm, kernel similarity, and $\Gamma(t)$.
\end{itemize}

\subsection{Representative PyTorch Snippet}

Below is a streamlined version (non-executable in this LaTeX format) capturing the essence of the V3 implementation:

\begin{verbatim}
class PrimeCosmosV3(nn.Module):
    def __init__(self, primes=[2,3,5,7], sector_dim=4,
                 alpha=1.5, beta_kernel=2.0, epsilon0=0.3):
        super().__init__()
        self.primes = primes
        self.n_p = len(primes)
        self.sector_dim = sector_dim
        self.alpha = alpha
        self.beta_kernel = beta_kernel
        self.epsilon0 = epsilon0

        # Lambda_m: Universal Multiplicity Constant (effective scalar)
        self.Lambda_m = nn.Parameter(torch.tensor(0.08, dtype=torch.float))

        # Local prime-sector tensors (grammar-aware)
        self.local_tensors = nn.ModuleDict({
            str(p): nn.Linear(sector_dim, sector_dim,
                             bias=False, dtype=torch.cfloat)
            for p in primes
        })

        # Correlation and coherence matrices
        self.C = nn.Parameter(torch.randn(self.n_p, self.n_p,
                                          dtype=torch.cfloat) * 0.1)
        self.gamma = nn.Parameter(torch.randn(self.n_p, self.n_p,
                                              dtype=torch.cfloat) * 0.05)

        # Grammar mask pi(p,q)
        self.grammar_mask = torch.zeros(self.n_p, self.n_p,
                                        dtype=torch.bool)
        for i, p in enumerate(primes):
            for j, q in enumerate(primes):
                self.grammar_mask[i, j] = (
                    q % p == 0 or p % q == 0 or
                    abs(math.log(p * q)) < 6.0
                )

    def multiplicity_kernel(self, X, target):
        # X, target: (B, N, d), real part for exponents
        v = X.real[..., :self.n_p]     # exponent proxies
        v_t = target.real[..., :self.n_p]
        d_mult = torch.norm(v - v_t, dim=-1, p=2)**2
        return torch.exp(-self.beta_kernel * d_mult)

    def compute_Gamma(self):
        abs_C = torch.abs(self.C)
        abs_gamma = torch.abs(self.gamma)
        return (abs_C * abs_gamma).sum() / (abs_C.sum() + 1e-8)

    def forward(self, X_t, target, noise_level=0.2):
        # XI(t) X_t
        Xi_X = torch.zeros_like(X_t, dtype=torch.cfloat)
        device = X_t.device

        for i, p in enumerate(self.primes):
            w_p = float(p) ** (-self.alpha)
            gamma_pq = self.gamma[i, :]  # (N,)
            logs = torch.log(
                torch.tensor([float(p * q) for q in self.primes],
                             device=device)
            )
            num = (logs * gamma_pq).sum()
            den = torch.abs(gamma_pq).sum() + 1e-8
            theta_p = num / den
            phase = torch.exp(1j * theta_p).unsqueeze(-1)

            A_p = self.local_tensors[str(p)](X_t[:, i, :])
            mask = self.grammar_mask[i].unsqueeze(-1).float()
            A_p = A_p * mask
            Xi_X[:, i, :] = w_p * phase * A_p

        # T_det and T_stoch
        kernel_sim = self.multiplicity_kernel(X_t, target)
        T_det = -0.02 * (X_t - target) * kernel_sim.unsqueeze(-1) - 0.001 * X_t
        Gamma_t = self.compute_Gamma()
        T_stoch = 0.05 * Gamma_t * torch.randn_like(X_t,
                                                    dtype=torch.cfloat)
        T_X = T_det + T_stoch

        # Sector-wise Lambda_m^{op}(t)
        Lambda_op = self.Lambda_m * torch.exp(
            -0.1 * torch.arange(1, self.n_p + 1,
                                device=device, dtype=torch.float)
        )
        correction = Lambda_op.unsqueeze(0).unsqueeze(-1) * T_X

        X_tp1 = Xi_X + correction + noise_level * \
                torch.randn_like(X_t, dtype=torch.cfloat)

        metrics = {
            'norm': torch.norm(X_tp1).item(),
            'kernel_sim': kernel_sim.mean().item(),
            'Gamma': Gamma_t.item()
        }
        return X_tp1, metrics
\end{verbatim}

\section{Factorization Dataset and Experimental Protocol}

\subsection{Toy Factorization Dataset}

We consider primes $P_N=\{2,3,5,7\}$ and an exponent cap $E_{\max}=3$. For each integer $n$ representable as $n=\prod_{i=1}^N p_i^{e_i}$ with $e_i\le E_{\max}$, we define an encoding
\[
Y^{(n)}\in\mathbb{R}^{1\times N\times (E_{\max}+1)},
\]
where for each prime $p_i$ we set
\[
Y^{(n)}_{0,i,e_i} = 1, \quad Y^{(n)}_{0,i,e} = 0 \text{ for } e\neq e_i.
\]

We construct:
\begin{itemize}[leftmargin=1.5em]
\item A \textbf{training set} of ``binary composites'': numbers $n=p_i p_j$ for $p_i,p_j\in P_N$ (with replacement), i.e., all products of 2 primes (possibly equal).
\item A \textbf{test set} of unseen ``triples'' and higher-degree composites, e.g., $n=30=2\cdot 3\cdot 5$, $n=42=2\cdot 3\cdot 7$, numbers with exponents $>1$ such as $12=2^2\cdot 3$.
\end{itemize}

The loss function is defined via the multiplicity kernel. For a predicted state $X_T$ and target $Y^{(n)}$, we may use
\[
\mathcal{L}(X_T,Y^{(n)}) = 1 - \bar{k}_{\mathrm{mult}}(X_T,Y^{(n)}),
\]
where $\bar{k}_{\mathrm{mult}}$ is the batch-averaged multiplicity kernel applied to the real exponent-channels of $X_T$ and $Y^{(n)}$.

\subsection{Training Protocol}

\begin{enumerate}[leftmargin=1.5em]
\item Initialize model parameters (local tensors, $C$, $\gamma$, $\Lambda_m$ in the safe band) and an optimizer, e.g.\ Adam.
\item For each epoch:
\begin{itemize}[leftmargin=1em]
\item For each training example $(n, Y^{(n)})$:
\begin{itemize}[leftmargin=0.8em]
\item Sample a noisy initial state $X_0$.
\item Recurrently apply the PrimeCosmosV3 forward pass for $T$ steps:
\[
X_{t+1},\text{metrics}_t = \mathrm{model}(X_t,Y^{(n)}).
\]
\item Compute loss at the final step, backpropagate through the unrolled steps, and update parameters.
\end{itemize}
\end{itemize}
\item Monitor training loss and average kernel similarity on the training set.
\end{enumerate}

\subsection{Evaluation Protocol}

After training:

\begin{itemize}[leftmargin=1.5em]
\item For each unseen composite $n$ in the test set:
\begin{itemize}[leftmargin=0.8em]
\item Sample multiple initial states $X_0$ (noisy), apply $T$ recurrent steps with fixed parameters.
\item Record final kernel similarity, $\Gamma(T)$, and norms.
\end{itemize}
\item Aggregate results to obtain mean and variance of test similarity.
\end{itemize}

We then perform a $\Lambda_m$-sweep:
\begin{itemize}[leftmargin=1.5em]
\item For each candidate value of $\Lambda_m$ (within and slightly outside the safe band):
\begin{itemize}[leftmargin=0.8em]
\item Fix $\Lambda_m$ in the trained model (or fine-tune).
\item Run the same evaluation protocol.
\item Record: mean train similarity, mean test similarity, mean $\Gamma(T)$, and collapse/blow-up rates (based on norm thresholds).
\end{itemize}
\end{itemize}

The result is a mapping
\[
\Lambda_m \mapsto \Big(\text{train sim}, \text{test sim}, \Gamma_{\mathrm{final}},
\text{collapse rate}\Big),
\]
from which a critical band for $\Lambda_m$ can be identified.

\section{Novelty and Defensive Publication Scope}

\subsection{Novel Conceptual Integration}

The unique aspects of this architecture include:
\begin{itemize}[leftmargin=1.5em]
\item Treating memory as a \emph{prime-graded multiplicity space}, rather than as unstructured vectors or generic Hilbert spaces.
\item Using a \emph{universal constant} $\Lambda_m$ as a tunable scalar controlling the strength of nonlinear corrections, with provable stability conditions directly linking $\Lambda_m$ to an operator norm ratio $\varepsilon_0/L$.
\item Embedding the original NP-SUPERMEMORY notions---superposition, entanglement, tensor networks, stochastic coherence---into a \emph{single, dynamical law} with explicit operators $\Xi(t)$, $T$, $\Sigma(t)$, and correlation structures $C_{pq}$, $\gamma_{pq}(t)$.
\item Providing a concrete, trainable module and dataset with a detailed experimental protocol for probing coherence, critical surfaces, and compositional generalization.
\end{itemize}

\subsection{What This Document Covers as Prior Art}

This document is intended as a defensive publication for the following categories:

\begin{enumerate}[leftmargin=1.5em]
\item \textbf{Prime-labeled Hilbert modules and multiplicity kernels} applied to memory and learning systems.
\item \textbf{Dynamical laws of the form} $X_{t+1}=\Xi(t)X_t+\Lambda_m^{\mathrm{op}}T(X_t)+\Sigma\xi_t$ where:
\begin{itemize}[leftmargin=0.8em]
\item $\Xi(t)$ is constrained by a prime-based grammar for connectivity.
\item $T$ is a Lipschitz multiplicity-structured mapping using a multiplicity kernel over exponent vectors.
\item $\Lambda_m$ is a universal scalar controlling the strength of $T$ and thereby stability, with explicit Grönwall-type bounds.
\end{itemize}
\item \textbf{Trainable RNN architectures} that implement such dynamics for prime-factorization-like tasks and evaluate compositional generalization via a multiplicity kernel.
\item \textbf{Experimental protocols} for sweeping $\Lambda_m$, monitoring order parameters like $\Gamma(t)$, and identifying critical coherence surfaces in prime-graded superpositional memory systems.
\end{enumerate}

\section{Conclusion and Future Directions}

The developments documented here turn NP-SUPERMEMORY and the Universal Multiplicity Constant $\Lambda_m$ from conceptual entities into a mathematically grounded and experimentally testable architecture. The combination of:
\begin{itemize}[leftmargin=1.5em]
\item multiplicity-graded Hilbert modules,
\item Gaussian multiplicity kernels on exponent vectors,
\item $\Lambda_m$-regulated dynamical laws with proven mean-square stability,
\item and a concrete PrimeCosmos V3 PyTorch implementation
\end{itemize}
constitutes a distinctive prior-art position for prime-based, superpositional memory architectures in machine learning and cognitive modeling.[file:1]

Future work includes:
\begin{itemize}[leftmargin=1.5em]
\item Extending tasks beyond factorization to more general symbolic and combinatorial domains where prime-like multiplicities naturally arise.
\item Learning a schedule or function $\Lambda_m(t)$ instead of a constant, and characterizing the resulting dynamical regimes.
\item Analyzing continuous-time limits and deriving more explicit phase-transition equations for $\Gamma(t)$.
\end{itemize}

\appendix
\section*{Mathematical Appendix}
\addcontentsline{toc}{section}{Mathematical Appendix}

\section{Operator Norm Bounds and Stability}

In this appendix we formalize the operator norm bounds and stability results underlying the $\Lambda_m$-regulated dynamics described in the main text. We work in a general Hilbert space setting and then specialize to the prime-graded memory architecture.

\subsection{Setting and Notation}

Let $(\mathcal{H},\langle\cdot,\cdot\rangle)$ be a complex Hilbert space with induced norm $\|x\| = \sqrt{\langle x,x\rangle}$. We consider discrete-time dynamics
\begin{equation}
X_{t+1} = \Xi(t) X_t + \Lambda_m(t) T(X_t) + \Sigma(t)\,\xi_t,
\label{eq:dyn-general}
\end{equation}
where:
\begin{itemize}
\item $\Xi(t):\mathcal{H}\to\mathcal{H}$ is a bounded linear operator for each $t$.
\item $\Lambda_m(t)\in\mathbb{R}$ is a scalar multiplicity regulator, possibly time-varying but bounded.
\item $T:\mathcal{H}\to\mathcal{H}$ is a (globally) Lipschitz mapping with $T(0)=0$.
\item $\Sigma(t):\mathcal{H}\to\mathcal{H}$ is a bounded linear operator controlling noise injection.
\item $\xi_t$ is a random element (e.g.\ i.i.d.\ across $t$) with zero mean and bounded second moment $\mathbb{E}\|\xi_t\|^2\le \sigma_\xi^2$.
\end{itemize}

We assume the following uniform bounds:
\begin{align}
\|\Xi(t)\|_{\mathrm{op}} &:= \sup_{\|x\|=1} \|\Xi(t)x\| \le 1 - \varepsilon_0,\qquad &&\forall t,\label{eq:bound-xi}\\
\|\Sigma(t)\|_{\mathrm{op}} &\le \sigma_s,\qquad &&\forall t,\label{eq:bound-sigma}\\
\sup_t |\Lambda_m(t)| &\le \Lambda_{\max} < \infty,\label{eq:bound-lambda}
\end{align}
for some $\varepsilon_0\in(0,1)$, $\sigma_s\ge 0$, and $\Lambda_{\max}\ge 0$. We also assume:

\begin{equation}
\|T(x) - T(y)\|\le L\,\|x-y\|,\quad \forall x,y\in\mathcal{H},\qquad T(0)=0.
\label{eq:Lipschitz-T}
\end{equation}

\subsection{Deterministic Lipschitz Bound for $T$}

\begin{lemma}[Norm bound for $T$]
Under \eqref{eq:Lipschitz-T} and $T(0)=0$ we have
\[
\|T(x)\| \le L\,\|x\|,\qquad \forall x\in\mathcal{H}.
\]
\end{lemma}

\begin{proof}
Setting $y=0$ in \eqref{eq:Lipschitz-T} yields
\[
\|T(x) - T(0)\| \le L\,\|x-0\| \Rightarrow \|T(x)\|\le L\,\|x\|.
\]
\end{proof}

In the NP-SUPERMEMORY context, $T_{\det}$ is constructed as a composition of bounded linear maps and bounded pointwise nonlinearities (e.g.\ multiplicity kernels, damping). Because all components are Lipschitz with finite constants, the composite $T_{\det}$ is Lipschitz with some finite $L_{\det}$; similarly for $T$ when we include or bound $T_{\stoch}$ in expectation. For the deterministically analyzed part, we take $T=T_{\det}$.

\subsection{One-Step Mean Square Bound}

Define the deterministic part
\[
f_t(x) = \Xi(t)x + \Lambda_m(t)T(x),
\]
and the noise term
\[
\eta_t = \Sigma(t)\xi_t,
\]
so the dynamics \eqref{eq:dyn-general} become
\[
X_{t+1} = f_t(X_t) + \eta_t.
\]

\begin{lemma}[One-step bound]
Assume \eqref{eq:bound-xi}--\eqref{eq:Lipschitz-T} and define
\[
c_t := |\Lambda_m(t)| L,\qquad c := \sup_t c_t.
\]
Then for all $t$,
\[
\|f_t(x)\| \le (1-\varepsilon_0 + c)\,\|x\|.
\]
\end{lemma}

\begin{proof}
Using the triangle inequality and the operator norms:
\begin{align*}
\|f_t(x)\|
&= \|\Xi(t)x + \Lambda_m(t)T(x)\| \\
&\le \|\Xi(t)x\| + |\Lambda_m(t)|\,\|T(x)\| \\
&\le (1-\varepsilon_0)\,\|x\| + |\Lambda_m(t)|\,L\,\|x\| \\
&= (1-\varepsilon_0 + c_t)\,\|x\| \\
&\le (1-\varepsilon_0 + c)\,\|x\|.
\end{align*}
\end{proof}

\begin{corollary}[Mean square estimate]
If $c<\varepsilon_0$, i.e.\ $\sup_t |\Lambda_m(t)| L < \varepsilon_0$, then with $q:=(1-\varepsilon_0 + c)^2<1$ we obtain
\[
\|f_t(x)\|^2 \le q\,\|x\|^2.
\]
\end{corollary}

\begin{proof}
Immediate from the previous lemma:
\[
\|f_t(x)\|^2 \le (1-\varepsilon_0 + c)^2\,\|x\|^2 = q\,\|x\|^2,
\]
and $q<1$ because $1-\varepsilon_0 + c < 1$.
\end{proof}

\subsection{Incorporating Noise}

Assume $\xi_t$ is independent of $X_t$ and has $\mathbb{E}[\xi_t]=0$ and $\mathbb{E}\|\xi_t\|^2\le \sigma_\xi^2$. Then
\[
\eta_t = \Sigma(t)\xi_t,\qquad \mathbb{E}[\eta_t]=0,\qquad 
\mathbb{E}\|\eta_t\|^2 \le \sigma_s^2 \sigma_\xi^2.
\]

\begin{lemma}[Mean-square recurrence]
Under the above assumptions,
\[
\mathbb{E}\|X_{t+1}\|^2 \le q\,\mathbb{E}\|X_t\|^2 + C,
\]
where $q=(1-\varepsilon_0 + c)^2<1$ and $C := \sigma_s^2 \sigma_\xi^2$.
\end{lemma}

\begin{proof}
We have
\[
X_{t+1} = f_t(X_t) + \eta_t.
\]
Using the inequality $\|a+b\|^2\le 2\|a\|^2 + 2\|b\|^2$,
\[
\|X_{t+1}\|^2 \le 2\|f_t(X_t)\|^2 + 2\|\eta_t\|^2.
\]
Taking expectations and applying the mean square bound on $f_t$ and $\eta_t$:
\begin{align*}
\mathbb{E}\|X_{t+1}\|^2
&\le 2\,\mathbb{E}\|f_t(X_t)\|^2 + 2\,\mathbb{E}\|\eta_t\|^2 \\
&\le 2q\,\mathbb{E}\|X_t\|^2 + 2\sigma_s^2\sigma_\xi^2.
\end{align*}
Renaming $2q\to q$ and $2\sigma_s^2\sigma_\xi^2\to C$ (or just absorbing the factor 2 into $q$ and $C$ for looser bounds) yields
\[
\mathbb{E}\|X_{t+1}\|^2 \le q\,\mathbb{E}\|X_t\|^2 + C,
\]
with $q<1$.
\end{proof}

\subsection{Uniform Mean-Square Boundedness}

Consider the recurrence $a_{t+1}\le q a_t + C$ with $q<1$. This unrolls to
\[
a_t \le q^t a_0 + \frac{C}{1-q}\big(1-q^t\big)\le q^t a_0 + \frac{C}{1-q}.
\]

\begin{theorem}[Uniform mean-square boundedness]
If $q<1$ as above, then
\[
\limsup_{t\to\infty} \mathbb{E}\|X_t\|^2 \le \frac{C}{1-q} < \infty.
\]
\end{theorem}

\begin{proof}
Let $a_t := \mathbb{E}\|X_t\|^2$. From the previous lemma, $a_{t+1}\le q a_t + C$ with $0\le q<1$. Induction yields
\[
a_t \le q^t a_0 + \frac{C}{1-q}(1-q^t).
\]
Taking $\limsup$ on both sides as $t\to\infty$, the $q^t a_0$ term vanishes, and we get
\[
\limsup_{t\to\infty} a_t \le \frac{C}{1-q}.
\]
\end{proof}

\subsection{Finite-Horizon Almost-Sure Boundedness}

\begin{corollary}[Finite-horizon almost-sure bound]
For any fixed horizon $T\in\mathbb{N}$ and radius $R>0$,
\[
\mathbb{P}\big(\|X_t\| > R\big) \le \frac{\mathbb{E}\|X_t\|^2}{R^2},\qquad \forall t\le T.
\]
Thus, with $R$ chosen large enough, $X_t$ is almost surely bounded on $t\le T$.
\end{corollary}

\begin{proof}
This is Markov's inequality applied to the nonnegative random variable $\|X_t\|^2$:
\[
\mathbb{P}(\|X_t\| > R) = \mathbb{P}(\|X_t\|^2 > R^2) \le \frac{\mathbb{E}\|X_t\|^2}{R^2}.
\]
Because $\mathbb{E}\|X_t\|^2$ is finite for each $t$, we can choose $R$ such that the right-hand side is arbitrarily small.
\end{proof}

\section{Specialization to Prime-Graded Superpositional Memory}

We now connect the abstract bounds to the prime-graded NP-SUPERMEMORY setting.

\subsection{Prime-Sector Decomposition and Operator Norms}

In the NP-SUPERMEMORY implementation:
\begin{itemize}
\item $\mathcal{H}$ is instantiated as $\mathbb{C}^{N\times d}$ (or a subspace), where $N$ is the number of primes and $d$ the sector dimension.
\item $\Xi(t)$ acts as a block-diagonal operator with blocks $w_p(t) U_p(t)$, where $w_p(t)$ are scalar weights (e.g.\ $p^{-\alpha}$) and $U_p(t)$ are linear maps defined via local tensors and phases.
\end{itemize}

Hence,
\[
\|\Xi(t)\|_{\mathrm{op}} = \max_{p} \big(|w_p(t)|\,\|U_p(t)\|_{\mathrm{op}}\big).
\]
In practice, $U_p(t)$ are small complex linear layers with bounded spectral norms, and $w_p(t)$ are deterministic functions of primes and hyperparameters. By construction we choose or rescale so that
\[
\max_{p}|w_p(t)|\,\|U_p(t)\|_{\mathrm{op}} \le 1 - \varepsilon_0
\]
for some $\varepsilon_0>0$, ensuring the first assumption \eqref{eq:bound-xi} holds.

Similarly, the deterministic component $T_{\det}$ is constructed as
\[
T_{\det}(X) = -\eta (X - X_{\mathrm{target}})\odot k_{\mathrm{mult}}(X,X_{\mathrm{target}}) - \mu X + (\text{coherence centering}),
\]
where all operations are Lipschitz, dominated by a linear term plus a bounded kernel modulation. It follows that $T_{\det}$ is globally Lipschitz with some constant $L_{\det}$ that can be estimated numerically: e.g.\ by sampling random $X$ in a neighborhood of 0 and computing
\[
\hat{L}_{\det} \approx \sup_{i} \frac{\|T_{\det}(X_i)\|}{\|X_i\|}.
\]
Provided $\Lambda_m L_{\det}<\varepsilon_0$, the deterministic component of the dynamics obeys the contraction assumptions.

\subsection{Safe Band for $\Lambda_m$ in Practice}

Combining the general theory with the prime-graded realization yields:

\begin{theorem}[Safe band for $\Lambda_m$ in PrimeCosmos]
Let $\varepsilon_0>0$ satisfy $\sup_t\|\Xi(t)\|\le 1-\varepsilon_0$, and let $L_{\det}$ be a Lipschitz constant for $T_{\det}$ in the NP-SUPERMEMORY architecture. If a scalar $\Lambda_m$ satisfies
\[
|\Lambda_m| < \frac{\varepsilon_0}{L_{\det}},
\]
then the dynamics
\[
X_{t+1} = \Xi(t) X_t + \Lambda_m T_{\det}(X_t) + \Sigma(t)\xi_t
\]
are uniformly mean-square bounded with
\[
\limsup_{t\to\infty}\mathbb{E}\|X_t\|^2 \le \frac{C}{1-q},
\]
for some finite $C$ and $q<1$ depending on $\varepsilon_0,\Lambda_m,L_{\det},\Sigma$.
\end{theorem}

This theorem captures the essence of the ``safe band'' interpretation: $\Lambda_m$ cannot be arbitrarily large if we require stable superpositional memory; its magnitude is constrained by the operator norm of $\Xi(t)$ and the Lipschitz constant of the multiplicity-preserving correction $T_{\det}$.

\subsection{Order Parameter and Phase Transition Heuristics}

While the above results provide rigorous mean-square bounds, the behavior of the coherence order parameter
\[
\Gamma(t) = \frac{\sum_{p,q}|C_{pq}|\,|\gamma_{pq}(t)|}{\sum_{p,q}|C_{pq}| + 1e-8}
\]
under variation of $\Lambda_m$ is analyzed empirically in experiments. The critical surface concept is:

\begin{itemize}
\item For $\Lambda_m$ within the safe band and moderate noise, empirical trajectories show $\Gamma(t)$ converging to a nonzero plateau, indicating stable coherence.
\item For $\Lambda_m$ exceeding a certain threshold (still possibly within the mathematical bound if noise or modeling approximations are changed), trajectories may exhibit either rapid decoherence ($\Gamma(t)\to 0$) or unstable behavior, suggesting empirical breakdown of effective memory even before rigorous mean-square instability manifests.
\end{itemize}

In this sense, the Grönwall-based safe band provides a necessary condition for stability, while the empirical critical surface (observed via sweeps of $\Lambda_m$, measures of $\Gamma(t)$, and compositional performance) characterizes where the system transitions between qualitatively different coherence regimes.

\section{Conclusion of the Appendix}

We have provided explicit operator norm bounds and mean-square stability results for the $\Lambda_m$-regulated dynamics used in NP-SUPERMEMORY, together with their specialization to the prime-graded superpositional memory architecture. The key point is that the Universal Multiplicity Constant $\Lambda_m$ is mathematically constrained by the interplay of $\Xi(t)$ and $T_{\det}$, rather than being a free parameter: its safe magnitudes must obey a ratio constraint $\Lambda_m L_{\det}<\varepsilon_0$ to guarantee bounded superpositional memory evolution.

\nocite{*}
\bibliographystyle{unsrt}  
\bibliography{references}  

\end{document}
